\citet{baptiste2008lagrangian} introduced the JITJSS problem and established the standard 72 benchmark test instances.
Since then, researchers have developed various exact, heuristic, and hybrid approaches.
Initial exact and hybrid approaches combined Lagrangian relaxation, branch-and-bound, and local search \citep{baptiste2008lagrangian, monette2009just}, as well as self-adapting Large Neighborhood Search \citep{laborie2007self}.

Evolutionary algorithms introduced effective sequence-level search strategies for JITJSS.
\citet{araujo2009genetic} and \citet{dos2010combination} proposed genetic algorithms, with \citet{dos2010combination} introducing a two-stage decomposition that pairs metaheuristic sequence search with Linear Programming to determine optimal operation start times.
\citet{yang2012improved} further expanded evolutionary methods by introducing a three-stage chromosome decoding strategy.

Trajectory methods, particularly Variable Neighborhood Search (VNS), established strong baselines by embedding exact sub-solvers into local search.
\citet{wang2014variable} combined VNS with linear programming schedule generation, while \citet{ahmadian2021meta} extended VNS by embedding integer programming directly within disjunctive two-machine local search operators, establishing the primary benchmark baseline across all 72 instances.

More recently, adaptive and learning-based techniques have been proposed.
\citet{abderrazzak2022adaptive} introduced an Adaptive Large Neighborhood Search (ALNS) driven by destroy-and-reconstruct operators, while \citet{sabri2023reinforcement} proposed a Deep Reinforcement Learning framework for rapid schedule synthesis.

\citet{wang2014variable} introduced effective penalty-guided neighborhoods within VNS, successfully addressing local trade-offs between earliness and tardiness within contiguous machine sequences.
While these operators optimize local alignments effectively, they operate independently of the multi-machine disjunctive critical paths that propagate delays across jobs.
Our proposed \texttt{swap\_crit} operator acts in a complementary manner at this structural level, tracing individual backward critical paths to target the broader predecessor chains that produce deadline violations.

% Overall, five key studies in the literature hold the best-known solution value for at least one benchmark instance: \citet{dos2010combination}, \citet{wang2014variable}, \citet{ahmadian2021meta}, \citet{abderrazzak2022adaptive}, and \citet{sabri2023reinforcement}.
% The composite literature benchmark for JITJSS is therefore defined by combining the lowest solution values achieved across these five works.
